% !TEX root = ../tecnologia-en-marcha.tex
% !TEX encoding = UTF-8 Unicode
% Texto: primer borrador del artículo, con acentos restaurados y citas enlazadas.
% Reordenado para abrir con la brecha de cumplimiento y presentar el monitoreo
% de personas adultas mayores como dominio de aplicación: el alcance de la
% investigación es la batería de pruebas, no el cuidado (PIA02, 1.1 y 1.7.2).

\section{Introducción}

La Estrategia Nacional de Inteligencia Artificial de Costa Rica plantea
principios rectores para el desarrollo y uso responsable de sistemas de IA,
entre ellos transparencia y explicabilidad, equidad y no discriminación, y
responsabilidad \cite{micitt2024enia}. Estos principios son coherentes con
instrumentos internacionales como el Reglamento Europeo de Inteligencia
Artificial, que establece obligaciones de transparencia, gobernanza de datos y
conservación de registros para sistemas de alto riesgo \cite{ue2024aiact}, y con
el NIST AI Risk Management Framework, que organiza la gestión de riesgos en
funciones de mapear, medir, gobernar y gestionar \cite{nist2023airmf}. El
Reglamento (UE) 2026/1744 aplazó la aplicación de las obligaciones de alto riesgo:
rigen desde el 2 de diciembre de 2027 para los sistemas del anexo III y desde el
2 de agosto de 2028 para los cubiertos por la legislación de producto del anexo I,
como la de dispositivos médicos \cite{ue2024aiact}. No obstante, estos marcos suelen expresarse en un nivel normativo o metodológico
que no siempre indica cómo producir evidencia técnica verificable sobre un
sistema específico.

El dominio de aplicación de este estudio son los sistemas que reconocen
actividades de la vida diaria a partir de sensores ambientales instalados en
viviendas, sin recurrir a cámaras ni micrófonos
\cite{kim2022inhome,requena2024independence}. En particular, los sensores de
infrarrojo pasivo (PIR), junto con sensores magnéticos de puerta y sensores de
temperatura, permiten registrar patrones de movimiento compatibles con tareas
como dormir, cocinar, comer o desplazarse dentro del hogar. Esta clase de
sistemas se ha propuesto para el monitoreo no intrusivo de personas adultas
mayores \cite{park2025aging}, una población para la que la propia ENIA reconoce
un potencial significativo de la IA \cite[p.~27]{micitt2024enia}. Sin embargo,
cuando las inferencias automatizadas se usan sobre poblaciones vulnerables, el
desempeño técnico del clasificador no es suficiente para justificar su uso.
% Las dos afirmaciones siguientes resumen la lectura que hace PIA02 (1.9 y
% 2.3.2) de esas revisiones.
Las revisiones recientes señalan, además, que la mayoría de los estudios del
campo se concentran en demostrar el desempeño técnico de los modelos y no en
evaluar su transparencia, su equidad o su gobernanza
\cite{park2025aging,habibovic2025aging}.

Otros trabajos proponen procedimientos completos para auditar o evaluar sistemas
de IA. Raji et al. definen un marco de auditoría algorítmica interna que acompaña
todo el ciclo de desarrollo y produce documentos en cada etapa
\cite{raji2020closing}. Zicari et al. proponen Z-Inspection, un proceso general de
inspección, basado en ética aplicada, para evaluar si un sistema cumple la
definición de IA confiable del grupo de expertos de la Comisión Europea
\cite{zicari2021zinspection}. Floridi et al. presentan capAI, un procedimiento de
evaluación de la conformidad con la propuesta de Reglamento Europeo de IA que
traduce principios éticos en criterios verificables, con un protocolo de revisión
interna y una ficha de evaluación externa \cite{floridi2022capai}. Estos marcos son
transversales a dominios y se organizan como procesos de revisión y documentación.
El marco de este estudio comparte su objetivo, pero parte de la ENIA y lo concreta
en pruebas computacionales automatizadas para sistemas basados en sensores PIR, que
producen un expediente verificable por terceros.

El problema abordado en este estudio es la brecha entre principios regulatorios
de alto nivel y mecanismos computacionales reproducibles para auditarlos en
sistemas de IA basados en sensores domiciliarios. En este dominio, la brecha es
particularmente relevante porque los datos describen rutinas domésticas
sensibles, los metadatos demográficos disponibles son limitados y los errores
pueden afectar de forma desigual a personas con patrones de movilidad o rutina
menos representados.

El objetivo de la investigación fue diseñar, implementar y verificar un marco
operativo que traduzca tres principios de la Estrategia Nacional de IA de Costa
Rica en requerimientos técnicos, pruebas computacionales y artefactos auditables
aplicables a un clasificador de actividades basado en sensores PIR. El objeto
de evaluación es el marco, no el sistema de monitoreo: el clasificador actúa
como sujeto de prueba.

Las principales contribuciones de este trabajo son cinco. Primero, una matriz de
nueve requerimientos verificables que traduce tres principios declarativos de la
ENIA en exigencias técnicas concretas, cada una con su evidencia esperada.
Segundo, un marco de cumplimiento implementado y probado, con tres módulos
operativos (explicabilidad, equidad y trazabilidad) y 292 pruebas automatizadas
que comprueban que el procedimiento es determinista y que un cambio posterior en
el protocolo queda en evidencia. Tercero, una validación empírica sobre dos
conjuntos de datos públicos con formatos distintos, sin modificar el núcleo del
marco, con 18\,736 inferencias registradas y un expediente completo y sellado por
hash en cada caso. Cuarto, evidencia creíble incluso cuando es desfavorable: el
veredicto de no cumplimiento, sostenido por un análisis de estabilidad con diez
semillas, muestra que el marco no es un sello automático de aprobación. Quinto,
una respuesta a un vacío geográfico y sectorial: el corpus revisado no permitió
identificar una guía sectorial que operacionalice estos principios en sistemas
basados en sensores PIR para el contexto costarricense, y su articulación con el
Reglamento Europeo de IA lo hace relevante de cara a esa regulación.

El artículo resume la arquitectura del marco, la metodología de implementación y
los resultados obtenidos en dos aplicaciones piloto sobre datos públicos del
proyecto CASAS.
